\section{A fully paid recursion with the new certificates}
\label{sec:schedule}

The local affine-plane certificate becomes a global improvement only after
the expensive evaluation of the cached group functions is included. We use the explicit $9\to13\to7$ schedule below to improve the
second-stage bottleneck as well as the root. This section states all additional
resource assumptions and checks the complete recurrence.

\subsection{Normalized inherited interface}

Let $s\ge2$, $\lambda=\lceil\log_2s\rceil$, and let the initial
bipartite Boolean tensor network have at most $s$ wires and maximum degree
at most $\lceil s^{9/10}\rceil$. Its entries lie in
$(\mathbb Z/2^{P_0}\mathbb Z)[V]/(V^\ell)$ with $P_0,\ell\le s$.
Let $\mathcal L$ be its description length. The permitted original entry
procedures take $(s+\mathcal L)^{O(1)}2^{O(s^{9/10})}$ time and polynomial
space, with fixed exponents and constants. The ring extension, canonical
creation-time lifts, exact deformation, and Hermite recovery are the
inherited normalized interface of the reviewed packets. Thus the result
below is a proved refinement of that interface; it does not independently
verify the upstream exact-counting algorithm.

At an active call with $M$ wires, group count $r$, and taper length $u$,
write
\[
 \mathcal K(M,u)=\kappa\frac{M}{\sqrt u},\qquad
 \kappa=\frac{\pi\sqrt2}{\ln2}.
\]
The inherited full-Cauchy entropy estimate says that if $h_{gh}$ is the
base-two logarithm of the exact list cardinality on directed group pair
$(g,h)$, then for every fixed nonnegative weight vector $z$,
\[
 \sum_{g\ne h}z_{gh}h_{gh}
 \le\left(\frac{\sqrt{\sum_{g\ne h}z_{gh}^2}}r+o(1)\right)
       \mathcal K(M,u).
\]
Exact weight allocations have total below the polynomially bounded working
precision. Because only the three fixed group counts $9,13,7$ occur, there
are polynomially many such allocations. The local certificates from the
preceding section apply uniformly to each allocation.

There are $2(r-1)$ directed pairs incident to one group. Consequently the
total number of singleton-table entries, including the sum over all groups,
is at most
\[
 \operatorname{poly}(s)\,
 2^{(\alpha_r+o(1))\mathcal K(M,u)},\qquad
 \alpha_r=\frac{\sqrt{2(r-1)}}r.
\]
Every entry is computed by a recursive child contraction with at most
$(2/r+o(1))M$ wires. Charging these calls separately from the local tree
operations gives
\begin{equation}
 T_d\le \operatorname{poly}(s)
 \left[2^{(\alpha_{r_d}+o(1))\mathcal K_d}T_{d+1}
       +2^{(b_{r_d}+o(1))\mathcal K_d}\right],
 \qquad \mathcal K_d=\mathcal K(M_d,u_d).
 \label{eq:affine-paid-recurrence}
\end{equation}
The local coefficients used here are
\[
 b_9^2=\frac13,\qquad b_{13}^2=\frac{601}{1690},\qquad
 b_7^2=\frac{632}{1225}.
\]
The first two come from the new four-plan certificates. The last is the
fixed seven-group plan $(0,1,2)\mid(3,4)\mid(5,6)$: its integer operation
supports have maximum squared norm $316$, so its directed normalized
coefficient is $2\cdot316/(25\cdot7^2)$. The supplied checker verifies this
bound directly.

The matrix operations use the inherited exact integer-coordinate and CRT
bridge for a fixed rational matrix multiplication algorithm with exponent
strictly below $12/5$. Thus the proof pays for the arithmetic over the
source ring rather than assuming that the ring is a field. No fast matrix
multiplication backend or end-to-end contraction solver is implemented by
the certificate programs.

\subsection{An explicit strict rational schedule}

Choose
\begin{align*}
 r_0&=9,& \eta_0&=\frac{15773}{1000000},\\
 r_1&=13,&\eta_1&=\frac9{800},\\
 r_d&=7\quad(d\ge2),&
 \eta_d&=\frac{8597}{1800000}
                   \left(\frac49\right)^{d-2}\quad(d\ge2).
\end{align*}
Put
\[
 u_d=\lfloor\eta_d\lambda\rfloor,\qquad
 B=\lceil6\log_2\lambda\rceil,\qquad
 a_d=\sum_{j=0}^d(u_j+B),\qquad
 H=\left\lceil\frac{s}{\lambda^{3/4}}\right\rceil.
\]
Use the inherited near-unit padding with relative slack $\lambda^{-3}$.
Calls with $M\le H$ are evaluated directly. For sufficiently large $s$
every active $u_d$ is positive; a fixed exceptional initial range is handled
by exhaustive contraction. The exact capacity sums are
\[
 \eta_0+\eta_1=\frac{27023}{1000000},\qquad
 \sum_{d\ge0}\eta_d=\frac{1781}{50000}.
\]

The following integer inequalities give rational upper bounds for every
logarithm needed in the packing proof:
\[
 9^{100000}<2^{316993},\qquad
 13^{100000}<2^{370044},\qquad
 7^{100000}<2^{280736}.
\]
They are checked using exact arbitrary-precision integers. Their respective
upper bounds for $\log_2 9$, $\log_2 13$, and $\log_2 7$ are
$316993/100000$, $370044/100000$, and $280736/100000$.

At the root both old sides have degree exponent $9/10$. At a later depth,
one side has that same old-degree bound, whereas the fresh side has degree
at most $2^{a_{d-1}+1}$. The asymmetric packing requirements are, up to
polylogarithmic factors, $Q^2\Delta$ at the root and $QR$, $Q^2C$,
$2^aQ^2$ later, with $Q=r^a$, retained degree $R\le s^{9/10+o(1)}$,
and fresh degree $C\le2^{a_{d-1}+1}$. The dominant margins are
\begin{center}
\begin{tabular}{lll}
\hline
Stage & Upper bound on the leading exponent & Margin below one\\
\hline
Root & $\displaystyle \frac9{10}
       +2\frac{15773}{1000000}\frac{316993}{100000}$
     & $\displaystyle\frac{69411}{50000000000}$\\[5pt]
Depth one & $\displaystyle \frac9{10}
       +\frac{27023}{1000000}\frac{370044}{100000}$
     & $\displaystyle\frac{75247}{25000000000}$\\[5pt]
Tail & $\displaystyle \frac9{10}
       +\frac{1781}{50000}\frac{280736}{100000}$
     & $\displaystyle\frac{287}{156250000}$\\
\hline
\end{tabular}
\end{center}
Every margin is strictly positive. The secondary fresh-side and packet
terms are controlled by
\[
 \frac{1781}{50000}
       \left(2\frac{370044}{100000}+1\right)<1.
\]
The fresh packet degree exponent is at most $1781/50000+o(1)<9/10$.
The smallest displayed power gaps are of order $10^{-6}$. They certify
eventual asymptotic feasibility, with no useful practical-size threshold
claimed here. This is one reason the coefficient comparison should not
be interpreted as a benchmark prediction.

\subsection{Uniformity, depth, and precision}

We give a bootstrap rather than using a depth bound as an unproved premise
of packing. Consider hypothetical active tail depths up to
$3\log_3\lambda/4+O(1)$. On this range,
\[
 u_d\ge\lambda^{\epsilon+o(1)},\qquad
 \epsilon=1-\frac{3\ln(9/4)}{4\ln3}>0.
\]
For example, positivity follows from $(9/4)^3<3^4$.
The cumulative guard adds only $O((\log\lambda)^2)$ to $a_d$.
The positive fixed margins above absorb this guard, all packing
polylogarithms, and $s/M<\lambda^{3/4}$. Padding is therefore valid
throughout the hypothetical range, with relative error
$O(\lambda^{-3})+s^{-c+o(1)}$ for a fixed $c>0$. In the tail every
child has at most $(2/7+o(1))M<M/3$ wires eventually. Hence an active
tail cannot extend beyond the stated hypothetical range. This closes the
bootstrap and proves the actual depth is $O(\log\lambda)$.

All active taper lengths tend to infinity uniformly. Floors therefore
change their leading terms by $o(1)$ relatively. The accumulated padding
error over the path is $1+o(1)$, and consequently
\begin{equation}
 M_0\le s,\qquad
 M_1\le(1+o(1))\frac{2s}{9},\qquad
 M_d\le(1+o(1))\frac{4s}{117}
                   \left(\frac27\right)^{d-2}\quad(d\ge2).
 \label{eq:affine-wire-masses}
\end{equation}

For a retained valuation $i<u_d$, set $D=2^{a_d-i}$. Then
$D\ge2^{a_{d-1}+B+1}$. The inherited input precision of the call is at
most $O(s\lambda^2 2^{a_{d-1}})$. Relative to $MD/u_d$, it is bounded by
\[
 O\!\left(\lambda^{15/4}2^{-B}\right)=o(1).
\]
The ordinary Hermite evaluation loss is lower order as well. The inherited
sharp precision bound $p_t\le(6+o(1))MD/u_d$, and hence the weighted
entropy estimate, therefore hold uniformly over all active calls.
All precisions retain polynomial bit length, since
$a_d\le(1781/50000)\lambda+O((\log\lambda)^2)$.

Direct leaf contraction costs $2^{O(s/\lambda^{3/4})}$ times the entry
cost. Original entries cost $2^{O(s^{9/10})}$ times a fixed polynomial.
Fresh packet entries cost
$\exp(O(2^{a_d}\log s))=\exp(s^{1781/50000+o(1)})$.
These terms are all $2^{o(s/\sqrt\lambda)}$. Their immutable
creation-time descriptions are retained; later entry evaluation does not
recursively expand earlier contraction routines. The ancestor cache costs
multiplying the leaf work remain explicitly charged in
\eqref{eq:affine-paid-recurrence}.

\subsection{Exact evaluation of the recurrence}

Define the leading coefficients of $\mathcal K_d$ on the common scale
$s/\sqrt\lambda$ by
\[
 K_0^*=\frac{\kappa}{\sqrt{\eta_0}},\qquad
 K_1^*=\frac{2\kappa}{9\sqrt{\eta_1}},\qquad
 K_2^*=\frac{4\kappa}{117\sqrt{\eta_2}}.
\]
After depth two the ratio is
\[
 \gamma=\frac{2/7}{\sqrt{4/9}}=\frac37.
\]
Unrolling the paid recurrence gives the maximum of three quantities:
\begin{align*}
 C_0&=b_9K_0^*,\\
 C_1&=\alpha_9K_0^*+b_{13}K_1^*,\\
 C_\infty&=\alpha_9K_0^*+\alpha_{13}K_1^*
              +K_2^*\max\left\{b_7,
                                  \frac{\alpha_7}{1-\gamma}\right\}.
\end{align*}
Indeed, the finite tail aggregation branches have tail part
\[
 K_2^*\left(\alpha_7\frac{1-\gamma^j}{1-\gamma}
                 +b_7\gamma^j\right),\qquad j\ge0,
\]
and the terminal leaf route is the limit. Here
$\alpha_7/(1-\gamma)=\sqrt3/2>b_7$, so the limit controls every tail
branch. This also pays for the route containing all recursive cache calls.

Exact radical expressions for these three coefficients are
\begin{align*}
 C_0&=\frac\pi{\ln2}\sqrt{\frac{2000000}{47319}},\\
 C_1&=\frac\pi{\ln2}\left(
      \sqrt{\frac{32000000}{1277613}}
       +\sqrt{\frac{384640}{123201}}\right),\\
 C_\infty&=\frac\pi{\ln2}\left(
      \sqrt{\frac{32000000}{1277613}}
       +\sqrt{\frac{51200}{41067}}
       +\sqrt{\frac{1600000}{4358679}}\right).
\end{align*}
The supplied exact-rational interval calculation gives
\begin{center}
\begin{tabular}{ll}
\hline
Branch & Coefficient\\
\hline
$C_0$ & $29.46603025659479065850\ldots$\\
$C_1$ & $30.69133782428594886802\ldots$\\
$C_\infty$ & $30.48972822403429948939\ldots$\\
\hline
\end{tabular}
\end{center}
In particular $C_1$ is the unique maximum of the three. The computation
encloses $\pi$ by Machin's identity and alternating arctangent series,
$\ln2$ by the positive series $2\operatorname{artanh}(1/3)$ with an explicit
tail bound, and square roots by integer square-root enclosures.
The interval endpoints are outward rounded; the table displays approximations.
Floating-point approximations are not used in any acceptance check.

For space, local tables are bounded by the chosen tree's local operation
bound, including singleton caches. Ancestor tables stored on one recursive
path add in size; their exponents do not add. Thus the space coefficient is
\[
 \max\{b_9K_0^*,\ b_{13}K_1^*,\ b_7K_2^*\}=C_0,
\]
since the latter two quantities are respectively
$8.0083771568264\ldots$ and $2.2775442472698\ldots$.
The fixed polynomial factors across $O(\log\lambda)$ levels contribute
only $O(\lambda\log\lambda)=o(s/\sqrt\lambda)$ to the exponent.
The dependence on the original description length remains one fixed
polynomial, because branching is independent of $\mathcal L$.

\begin{theorem}[New generic normalized certificate]
Under the inherited normalized interface specified above, the affine-plane
and four-plan construction contracts the root network deterministically in
\[
 (s+\mathcal L)^{O(1)}
 2^{(30.69133782428594886802\ldots+o(1))s/\sqrt\lambda}
\]
time and
\[
 (s+\mathcal L)^{O(1)}
 2^{(29.46603025659479065850\ldots+o(1))s/\sqrt\lambda}
\]
space. The eventual coefficients $30.692$ and $29.467$, respectively, are
valid. The method uses subexponential space.
\end{theorem}

\begin{proof}
The local selectors merely reassociate exact finite sums and products in
the inherited mode-contraction identity. Every feasible exact-weight
allocation is evaluated once. The strict packing, depth, precision, and
entry-access checks above justify all inherited calls. The paid recurrence
and local storage calculation give the asserted resources, including leaves
and polynomial overhead.
\end{proof}

Compared with the reviewed generic cached certificate
$30.891018920366\ldots$ in time and $30.753912832055\ldots$ in space,
both leading coefficients decrease. This is a comparison of explicit upper
bounds per normalized wire; it is neither a measured SAT speedup nor a
coefficient per SAT variable. The logarithmic order
$s/\sqrt{\log s}$, and consequently the inherited unrestricted SAT order,
remain unchanged. No global optimality of the finite portfolios or schedule
is claimed.
